\appendix
\section{Appendix: Corpus Design}

The corpus in this package is intentionally small and auditable. Each record stores:

\begin{itemize}[leftmargin=1.4em]
  \item date and location;
  \item author and recipient;
  \item phase of Wilkinson's career;
  \item a relationship-mode label;
  \item optional Frederick-role coding;
  \item a short excerpt and summary;
  \item source URL;
  \item named connections used for the network figure.
\end{itemize}

This design makes the paper easy to inspect and criticize. A reader can disagree with the coding while still seeing exactly how the paper's figures were produced. That transparency matters for a humanities workflow because it prevents the computational layer from turning into black-box authority.
